\clearpage\thispagestyle{empty}
\pdfbookmark[0]{Quellen zum Ceres-Bericht vom Dezember 1801}{ceres180112-sources}
\begin{center}\textbf{Quellen zum Ceres-Bericht vom Dezember 1801}\end{center}
\noindent Diese KI-Übertragung in heutiges Deutsch folgt dem deutschen Text in Carl Friedrich Gauß’ \href{https://gdz.sub.uni-goettingen.de/id/PPN236007467}{\textit{Werken}, Band~VI} (Göttingen, 1874), S.~199--204. Franz Xaver von Zach berichtet über Gauß’ Berechnungen und Bemerkungen.
\par\smallskip\noindent Von Zachs \href{https://archive.org/download/bub_gb_8mY7AQAAIAAJ/page/n544.jpg}{ursprünglicher Aufsatz}, Nr.~LVII, erschien im Dezemberheft 1801 der Zeitschrift \textit{Monatliche Correspondenz zur Beförderung der Erd- und Himmels-Kunde} (Gotha, Becker), Band~IV, S.~638--649. Der in den \textit{Werken} wiedergegebene Text setzt beim \href{https://archive.org/download/bub_gb_8mY7AQAAIAAJ/page/n545.jpg}{Absatz auf S.~639} ein und reicht bis zum \href{https://archive.org/download/bub_gb_8mY7AQAAIAAJ/page/n555.jpg}{Schluss auf S.~649}. Die Überschrift in den \textit{Werken} nennt Band~III. Die Digitalisate der Zeitschrift stammen vom Exemplar der University of California, das Google digitalisiert hat und das im Internet Archive zugänglich ist.
\par\smallskip\noindent Übertragung und Satz wurden mit KI erstellt.
\clearpage
